% ============================================================
%  Preambolo aggiuntivo per la dispensa di Programmazione
%  Estende unina_doc_class.cls (stile ripreso dalle Dispense APA e
%  dalla dispensa di Tecnologie Web, di cui riusa listati e riquadri)
% ============================================================

% Tabelle con i filetti di booktabs (\toprule, \midrule, \bottomrule)
\usepackage{booktabs}

% ─── Font aggiuntivi ─────────────────────────────────────────
% \Mononoki e' richiamato dalla classe nelle macro \dir e \file
\newfontfamily\Mononoki{mononoki}[
  Path = fonts/, Extension = .ttf,
  UprightFont = *-Regular, BoldFont = *-Bold,
  ItalicFont  = *-Italic,  BoldItalicFont = *-BoldItalic]

% FiraCode SENZA legature di programmazione.
% La classe dichiara \FiraCode con Contextuals=Alternate: in un testo
% sul C questo e' dannoso, perche' "!=" verrebbe disegnato come "≠", "->"
% come una freccia e "<=" come "≤", cioe' non come cio' che lo studente
% deve digitare. Qui si ridefinisce \FiraCode senza legature.
\newfontfamily\FiraCodeNL{FiraCode}[
  Path = fonts/, Extension = .ttf,
  UprightFont = *-Regular, BoldFont = *-Bold,
  Ligatures = NoCommon, Contextuals = NoAlternate]
\let\FiraCode\FiraCodeNL

% Nessun "$" deve aprire la modalita' matematica dentro un listato
\lstset{mathescape=false}

% ─── Etichette richieste dalla classe per l'elenco dei listati ─
\gdef\documentlsttitle{Elenco dei listati}
\gdef\documentlstlabel{Listato}

% ─── Colori per il codice su sfondo chiaro (figure e riquadri) ─
\definecolor{codestring}{HTML}{1F7A5A}
\definecolor{codecomment}{HTML}{8A8A8A}
\definecolor{codekey}{HTML}{C2185B}

% ============================================================
%  ASPETTO DEI LISTATI: UNA FINESTRA DI VS CODE (TEMA DARK+)
% ============================================================
% Ogni listato e' disegnato come una piccola finestra dell'editor Visual
% Studio Code con il tema scuro: sfondo quasi nero, numeri di riga grigi a
% sinistra e, se il listato ha un nome, la barra delle schede in alto con
% la scheda del file aperto (icona del linguaggio, nome del file e la
% crocetta di chiusura). E' cio' che lo studente vede sul proprio schermo.
\definecolor{vscbg}{HTML}{1E1E1E}        % sfondo dell'editor
\definecolor{vsctabbar}{HTML}{252526}    % barra delle schede
\definecolor{vscfg}{HTML}{D4D4D4}        % testo normale, punteggiatura
\definecolor{vsclinenum}{HTML}{858585}   % numeri di riga
\definecolor{vscblue}{HTML}{569CD6}      % tipi base, modificatori, macro
\definecolor{vscpurple}{HTML}{C586C0}    % if, for, return, #include...
\definecolor{vscstring}{HTML}{CE9178}    % stringhe e caratteri
\definecolor{vsccomment}{HTML}{6A9955}   % commenti
\definecolor{vscfunc}{HTML}{DCDCAA}      % nomi di funzione
\definecolor{vscvar}{HTML}{9CDCFE}       % variabili e parametri
\definecolor{vsctype}{HTML}{4EC9B0}      % tipi definiti (typedef, struct)
\definecolor{vscnum}{HTML}{B5CEA8}       % numeri
\definecolor{vscaccent}{HTML}{0078D4}    % filetto della scheda attiva
\definecolor{vsctabtext}{HTML}{FFFFFF}   % nome del file nella scheda
\definecolor{vsctabdim}{HTML}{9D9D9D}    % crocetta di chiusura
% colori delle icone dei file (tema di icone Seti, quello predefinito)
\definecolor{vsciconorange}{HTML}{E37933}
\definecolor{vsciconblue}{HTML}{519ABA}
\definecolor{vsciconyellow}{HTML}{CBCB41}
\definecolor{vscicongreen}{HTML}{8DC149}
\definecolor{vsciconpurple}{HTML}{A074C4}
\definecolor{vsciconred}{HTML}{CC3E44}

% ============================================================
%  LINGUAGGI PER LISTINGS
% ============================================================
% Listings non conosce la grammatica del C: colora le parole che compaiono
% in elenchi fissi. Le classi di parole chiave riproducono le categorie del
% tema Dark+ di VS Code (con l'estensione C/C++ installata):
%   [1] blu       tipi base, modificatori, sizeof e macro (NULL, EOF, PI...)
%   [2] viola     controllo del flusso (if, for, return...) e direttive (#include)
%   [3] giallo    funzioni chiamate o definite negli esempi della dispensa
%   [4] verde     tipi definiti con typedef e nomi delle struct
% Tutti gli altri identificatori sono variabili e prendono l'azzurro chiaro.
\lstset{keywords={}}

\lstdefinelanguage{C99}{
  sensitive=true,
  alsoletter={\#},
  morecomment=[l]{//},
  morecomment=[s]{/*}{*/},
  morestring=[b]",
  morestring=[b]',
  keywords=[1]{
    int,char,float,double,void,short,long,signed,unsigned,struct,union,
    enum,typedef,const,static,extern,register,auto,volatile,sizeof,inline,
    restrict,_Bool,bool,true,false,
    % macro di libreria e degli esempi
    NULL,EOF,PI,CLASS_SIZE,RAND_MAX,INT_MAX,INT_MIN,CHAR_MAX,CHAR_MIN,
    LONG_MAX,UINT_MAX,SHRT_MAX,EXIT_SUCCESS,EXIT_FAILURE,FLT_EPSILON,
    DBL_EPSILON,LIMIT,MAXWORD,N,M,DIM,MAX,SQ,SQUARE,DEBUG,BUFSIZ,
    STUDENTE_H,SORT_H,AREE_H,BIG,PAUSA,SCHAR_MIN,SCHAR_MAX,UCHAR_MAX,
    SHRT_MIN,LONG_MIN,ULONG_MAX,FLT_MIN,FLT_MAX,FLT_DIG,DBL_MIN,DBL_MAX,
    DBL_DIG,INITIAL_SEED,MULTIPLIER,INCREMENT,MODULUS,FLOATING_MODULUS},
  keywords=[2]{
    if,else,for,while,do,switch,case,default,break,continue,return,goto,
    \#include,\#define,\#ifndef,\#ifdef,\#endif,\#if,\#else,\#elif,\#undef,
    \#pragma},
  keywords=[3]{
    main,printf,scanf,getchar,putchar,puts,gets,fgets,fputs,fprintf,
    fscanf,sprintf,sscanf,fopen,fclose,malloc,calloc,realloc,free,exit,
    abs,fabs,sqrt,pow,sin,cos,exp,log,floor,ceil,rand,srand,time,
    strlen,strcpy,strncpy,strcat,strncat,strcmp,strncmp,strchr,strstr,
    strcspn,strspn,strdup,memcpy,memset,isalpha,isdigit,isspace,isupper,
    islower,isalnum,ispunct,toupper,tolower,assert,feof,qsort,
    % funzioni degli esempi
    max,min,maxmin,prn_info,area,area_cerchio,area_quadrato,perimetro,
    fattoriale,factorial,potenza,power,swap,sum,somma,sum_array,media,
    stampa,stampa_array,leggi,leggi_array,fail,update,inc,incrementa,
    wrt_address,twice,all_add,compute_sum,f,g,h,sf,count,counter,
    random_seed,my_random,ordina,bubble_sort,selection_sort,bubble,
    sort_words,wrt_words,error_exit,word_cnt,my_strlen,my_strcpy,
    my_strcat,crea_matrice,libera_matrice,get_matrix_space,
    release_matrix_space,fill_matrix,wrt_matrix,
    nuovoNodo,inserisciTesta,inserisciCoda,inserisciPosizione,
    inserisciOrdinato,cerca,lunghezza,eliminaValore,eliminaPosizione,
    eliminaTesta,eliminaCoda,unisci,stampaLista,liberaLista,
    leggiStudenti,stampaStudenti,inserisci,rimuovi,inserisciNodo,
    rimuoviNodo,stampaLista,cancellaLista,cnt_char,capitalize,
    trial,mult,prova,converti,fahrenheit,celsius,cubo,quadrato,
    leggi_intero,conta,scambia,inverti,cerca_massimo,mcd,primo,
    stampa_matrice,somma_matrici,somma_righe,trasponi,setbuf,sleep,pausa,
    errore,nothing,absolute_value,probability,random},
  keywords=[4]{
    size_t,FILE,ptrdiff_t,int8_t,int16_t,int32_t,int64_t,uint8_t,uint32_t,
    uint64_t,carta,complex,student,date,department,address,employee_data,
    fruit,vegetable,card,Node,Nodo,studente,Studente,lista,Lista,punto,
    Punto,persona,Persona,matrice,elemento,elem,frazione,data,
    stack,queue,Coda,Pila,uppercase,inches,u_char,tipoInt},
}

% ─── Pseudocodice ────────────────────────────────────────────
% Le parole chiave (in maiuscolo) sono viola come le istruzioni di
% controllo del C; i testi fra virgolette sono stringhe.
\lstdefinelanguage{Pseudo}{
  sensitive=true,
  morestring=[b]",
  morecomment=[l]{//},
  keywords=[1]{BEGIN,END,LEGGI,STAMPA,SE,ALLORA,ALTRIMENTI,MENTRE,
               FINCHE,RIPETI,PER,DA,A,ESEGUI,FINE,RESTITUISCI},
  keywords=[2]{PI,AND,OR,NOT},
}

% ─── Comandi del terminale e script di shell ─────────────────
% I comandi sono gialli, come le funzioni (VS Code fa lo stesso).
\lstdefinelanguage{Shell}{
  sensitive=true,
  morecomment=[l]{\#},
  morestring=[b]",
  morestring=[b]',
  morekeywords={gcc,cc,clang,make,cd,ls,mkdir,cat,echo,cp,mv,rm,pwd,
                sudo,apt,wsl,code,chmod,touch,man,clear,gdb,diff,
                xcode-select,bash,sh,history,less,head,tail},
}

% ─── Makefile ────────────────────────────────────────────────
% Commenti verdi, variabili $(...) azzurre, obiettivi gialli. Il TAB
% obbligatorio prima dei comandi si scrive (*@\TAB@*) ed e' disegnato
% come una freccia, come fa VS Code mostrando gli spazi bianchi.
\lstdefinelanguage{Make}{
  sensitive=true,
  alsoletter={.-},
  morecomment=[l]{\#},
  morestring=[b]",
  moredelim=**[s][\color{vscvar}]{\$(}{)},
  keywords=[1]{CC,CFLAGS,OBJ,OBJS,LDLIBS,LDFLAGS,TARGET,SRC},
  keywords=[2]{gcc,rm,echo,cc},
}

% ============================================================
%  STILI DI LISTATO
% ============================================================
% Ogni stile sceglie anche l'icona mostrata nella scheda (chiave vscicon,
% definita piu' sotto).
% I nomi delle intestazioni standard dopo #include (<stdio.h>...) sono
% stringhe per VS Code: li si colora con delle sostituzioni ("literate"),
% che con l'asterisco non agiscono dentro commenti e stringhe.
\lstdefinestyle{c}{
  language=C99,
  vscicon=c,
  keywordstyle=[1]\color{vscblue},
  keywordstyle=[2]\color{vscpurple},
  keywordstyle=[3]\color{vscfunc},
  keywordstyle=[4]\color{vsctype},
  identifierstyle=\color{vscvar},
  stringstyle=\color{vscstring},
  commentstyle=\color{vsccomment},
  literate=*%
    {<stdio.h>}{{\textcolor{vscstring}{<stdio.h>}}}9
    {<stdlib.h>}{{\textcolor{vscstring}{<stdlib.h>}}}{10}
    {<math.h>}{{\textcolor{vscstring}{<math.h>}}}8
    {<string.h>}{{\textcolor{vscstring}{<string.h>}}}{10}
    {<ctype.h>}{{\textcolor{vscstring}{<ctype.h>}}}9
    {<limits.h>}{{\textcolor{vscstring}{<limits.h>}}}{10}
    {<float.h>}{{\textcolor{vscstring}{<float.h>}}}9
    {<time.h>}{{\textcolor{vscstring}{<time.h>}}}8
    {<stdbool.h>}{{\textcolor{vscstring}{<stdbool.h>}}}{11}
    {<assert.h>}{{\textcolor{vscstring}{<assert.h>}}}{10}
    {<unistd.h>}{{\textcolor{vscstring}{<unistd.h>}}}{10},
}

% file di intestazione (.h): stessi colori, icona diversa
\lstdefinestyle{h}{
  style=c,
  vscicon=h,
}

\lstdefinestyle{shell}{
  language=Shell,
  vscicon=shell,
  keywordstyle=\color{vscfunc},
  stringstyle=\color{vscstring},
  commentstyle=\color{vsccomment},
}

\lstdefinestyle{make}{
  language=Make,
  vscicon=make,
  keywordstyle=[1]\color{vscvar},
  keywordstyle=[2]\color{vscfunc},
  stringstyle=\color{vscstring},
  commentstyle=\color{vsccomment},
}

\lstdefinestyle{plain}{
  language={},
  vscicon=plain,
}

\lstdefinestyle{pseudo}{
  language=Pseudo,
  vscicon=plain,
  keywordstyle=[1]\color{vscpurple},
  keywordstyle=[2]\color{vscblue},
  identifierstyle=\color{vscvar},
  stringstyle=\color{vscstring},
  commentstyle=\color{vsccomment},
}

% ============================================================
%  LA FINESTRA DELL'EDITOR
% ============================================================
\makeatletter
% Chiave vscicon: l'icona da mettere nella scheda (la sceglie lo stile)
\lst@Key{vscicon}{plain}{\def\wtvsc@icon{#1}}

% Il codice e' in corpo 7,5: la colonna dei numeri di riga toglie spazio
% orizzontale, e con il corpo leggermente ridotto le righe stanno in una
% riga anche nelle colonne strette di "affianco".
\newcommand{\wtvsccodefont}{\fontsize{7.5}{9}\selectfont\FiraCodeNL\color{vscfg}}
\newcommand{\wtvsccodefontsmall}{\fontsize{6.6}{7.6}\selectfont\FiraCodeNL\color{vscfg}}
\lstset{
  basicstyle=\linespread{1.15}\wtvsccodefont,
  % literate={} annulla le sostituzioni dello stile FiraCodeStyle (pacchetto
  % lstfiracode, caricato dalla classe): servono a tenere unite le coppie di
  % caratteri che Fira Code disegna come legature, ma cosi' listings non
  % riconosce piu' i delimitatori che le contengono (per esempio "//" e "/*"
  % dei commenti). Le legature sono comunque disattivate nel font.
  literate={},
  backgroundcolor=\color{vscbg},
  identifierstyle={},
  % numeri di riga, come nella "gutter" di VS Code
  numbers=left,
  numberstyle=\fontsize{5.5}{6}\selectfont\FiraCodeNL\color{vsclinenum},
  numbersep=5pt,
  % ogni listato riparte da 1
  firstnumber=1,
  % la cornice e' invisibile (spessore nullo): serve solo a estendere lo
  % sfondo scuro attorno al codice e sotto i numeri di riga
  frame=single,
  framerule=0pt,
  rulecolor=\color{vscbg},
  % il "titolo" del listato diventa la barra delle schede, sopra il codice
  captionpos=t,
  abovecaptionskip=0pt,
  belowcaptionskip=0pt,
  % un TAB vale 4 spazi, come in VS Code
  tabsize=4,
}

% Margini della finestra: la finestra va da #1 dal bordo sinistro della
% riga fino a #2 dal bordo destro. Dentro, a sinistra ci sono 16pt per la
% colonna dei numeri di riga (12pt) e lo spazio prima del codice (4pt);
% a destra 4pt di spazio.
\newcommand{\wtvscL}{0pt}
\newcommand{\wtvscR}{0pt}
\newcommand{\wtvscmargini}[2]{%
  \renewcommand{\wtvscL}{#1}\renewcommand{\wtvscR}{#2}%
  \lstset{xleftmargin=\dimexpr#1+16pt\relax, framexleftmargin=12pt,
          framesep=4pt,
          xrightmargin=\dimexpr#2+4pt\relax, framexrightmargin=0pt}}
\wtvscmargini{0pt}{0pt}

% Icone dei file nella scheda (riprendono il tema di icone Seti di VS Code)
\newcommand{\wtvsc@badge}[2]{%
  {\FiraCodeNL\bfseries\fontsize{6.5}{7}\selectfont\color{#1}#2}}
\newcommand{\wtvsc@iconmark}{%
  \ifcsname wtvsc@icon@\wtvsc@icon\endcsname
    \csname wtvsc@icon@\wtvsc@icon\endcsname
  \else
    \wtvsc@icon@plain
  \fi}
\def\wtvsc@icon@c{\wtvsc@badge{vsciconblue}{C}}
\def\wtvsc@icon@h{\wtvsc@badge{vsciconpurple}{h}}
\def\wtvsc@icon@make{\wtvsc@badge{vsciconorange}{M}}
\def\wtvsc@icon@shell{\wtvsc@badge{vscicongreen}{>\_}}
\def\wtvsc@icon@plain{{\fontsize{7}{7}\selectfont\color{vsctabdim}\ifilecode}}

% La barra delle schede: grigio scuro con gli angoli superiori arrotondati,
% e a sinistra la scheda attiva (dello stesso colore dell'editor, con il
% filetto blu in alto) che contiene icona, nome del file e crocetta.
\newlength{\wtvsc@barw}
\newcommand{\wtvsc@tabbar}{%
  \setlength{\wtvsc@barw}{\dimexpr\linewidth-\wtvscL-\wtvscR\relax}%
  \par\noindent\hspace*{\wtvscL}%
  \begin{tikzpicture}[baseline=0pt]
    \useasboundingbox (0,0) rectangle (\wtvsc@barw,17pt);
    \clip[rounded corners=3pt] (0,-5pt) rectangle (\wtvsc@barw,17pt);
    \fill[vsctabbar] (0,0) rectangle (\wtvsc@barw,17pt);
    \node[anchor=south west, inner xsep=7pt, inner ysep=0pt,
          minimum height=17pt, fill=vscbg, text=vsctabtext,
          font=\Lato\fontsize{7.5}{9}\selectfont] (tab) at (0,0)
      {\wtvsc@iconmark\hspace{4pt}\lstname\hspace{9pt}%
       {\color{vsctabdim}\fontsize{8}{9}\selectfont\texttimes}};
    \fill[vscaccent] (tab.north west) rectangle ([yshift=-1pt]tab.north east);
  \end{tikzpicture}%
  \par\nointerlineskip}

% La barra delle schede compare solo se il listato ha un nome (chiave
% "name"): se manca, il titolo viene svuotato appena prima di disegnarlo.
\lstset{title={\lstname}, everydisplay={\wtvsc@solosenome}}
\newcommand{\wtvsc@solosenome}{%
  \ifx\lst@intname\@empty
    \let\lst@caption\@empty
  \else
    \let\lst@caption\relax
  \fi}
% Al posto della didascalia si disegna la barra delle schede. La ridefinizione
% va fatta a inizio documento perche' il pacchetto caption, caricato prima,
% sostituisce anch'esso \lst@maketitle in quel momento.
\AtBeginDocument{\def\lst@maketitle#1{\wtvsc@tabbar}}
% Nel nome del file listings sostituisce "-" con un trattino lungo: nella
% scheda deve restare il trattino vero, quello del nome del file.
\def\lst@filenamerpl{_\textunderscore $\textdollar}

% Il TAB obbligatorio dei makefile, reso visibile: una freccia larga quanto
% quattro colonne del listato (\lst@width), cioe' quanto il TAB stesso. Si
% scrive (*@\TAB@*) all'inizio della riga di comando.
\newcommand{\TAB}{\makebox[4\lst@width][l]{%
  \color{vsclinenum}\tikz[baseline=-2.6pt]\draw[-{Stealth[length=3pt,width=3pt]},
    line width=0.5pt] (0.4pt,0) -- (3.4*\lst@width,0);}}
\makeatother

% ============================================================
%  CORREZIONE DEI RIQUADRI DELLA CLASSE
% ============================================================
% Nella classe il titolo viene passato a tcolorbox come "title=#1": una
% virgola nel titolo verrebbe interpretata come separatore di chiavi.
% Racchiudendo il titolo in un gruppo il problema sparisce.
\renewcommand{\info}[2]{\begin{infobox}{{#1}}#2\end{infobox}}
\renewcommand{\warning}[2]{\begin{warningbox}{{#1}}#2\end{warningbox}}
\renewcommand{\error}[2]{\begin{errorbox}{{#1}}#2\end{errorbox}}

% ============================================================
%  INDICE: NUMERI DI PAGINA A TRE CIFRE
% ============================================================
% I numeri di pagina a tre cifre non entrano nella casella standard di
% 1,55em; la casella si allarga a una misura fissa, cosi' tutti i numeri
% finiscono sullo stesso margine destro.
\makeatletter
\renewcommand{\@pnumwidth}{22pt}
\makeatother

% ============================================================
%  NUMERAZIONE DI ESEMPI ED ESERCIZI
% ============================================================
% Nella classe i contatori sono agganciati a "section"; con secnumdepth=0
% il contatore di sezione non viene mai incrementato, per cui si otterrebbe
% "Esempio 1.0.3". Li si riaggancia quindi al capitolo.
\counterwithout{esempio}{section}
\counterwithin{esempio}{chapter}
\counterwithout{esercizio}{section}
\counterwithin{esercizio}{chapter}

% L'esempio viene composto in una scatola: se sta in una pagina la si mette
% intera (e passa alla pagina dopo se non entra), altrimenti la si "svuota"
% nella pagina e l'esempio puo' spezzarsi normalmente.
\newbox\wtesbox
\newcommand{\wtemetti}{%
  \ifdim\dimexpr\ht\wtesbox+\dp\wtesbox\relax>\textheight
    \unvbox\wtesbox
  \else
    \box\wtesbox
  \fi}
\renewenvironment{esempio}[1]{%
  \par\addvspace{0.8em}%
  \setbox\wtesbox\vbox\bgroup
  \refstepcounter{esempio}%
  \noindent\rule{\linewidth}{0.5pt}\par\nopagebreak%
  \vspace{0.15em}%
  \noindent{\color{primary!80!black}\textbf{Esempio \theesempio}}%
  \ifx\relax#1\relax\else%
    {\color{primary!80!black}\textbf{ --- #1}}%
  \fi\par\nopagebreak%
  \vspace{0.3em}%
}{%
  \par\vspace{0.3em}%
  \noindent\rule{\linewidth}{0.5pt}\par
  \egroup
  \wtemetti
  \par\vspace{0.8em}%
}

\renewenvironment{esercizio}[1]{%
  \par\addvspace{0.8em}%
  \setbox\wtesbox\vbox\bgroup
  \refstepcounter{esercizio}%
  \noindent\rule{\linewidth}{1.2pt}\vspace{1pt}\par%
  \noindent\rule{\linewidth}{0.4pt}\par%
  \vspace{0.15em}%
  \noindent{\color{accent!80!black}\textbf{Esercizio \theesercizio}}%
  \ifx\relax#1\relax\else%
    {\color{accent!80!black}\textbf{ --- #1}}%
  \fi\par%
  \vspace{0.3em}%
}{%
  \par\vspace{0.3em}%
  \noindent\rule{\linewidth}{0.4pt}\vspace{1pt}\par%
  \noindent\rule{\linewidth}{1.2pt}\par
  \egroup
  \wtemetti
  \par\vspace{0.8em}%
}

% ============================================================
%  MACRO INLINE
% ============================================================
% Codice inline breve (non spezzabile). \leavevmode serve quando il codice
% apre un paragrafo (per esempio all'inizio della cella di una tabella):
% senza, il colore verrebbe applicato prima del paragrafo e la cella
% comincerebbe con una riga vuota.
\DeclareRobustCommand{\code}[1]{\leavevmode{\FiraCode\color{primary!65!black}\mbox{#1}}}
% Codice inline lungo (spezzabile a fine riga)
\DeclareRobustCommand{\codel}[1]{\leavevmode{\FiraCode\color{primary!65!black}#1}}
% Termine tecnico evidenziato (spezzabile, a differenza di \hi)
\def\term#1{{\color{accent!80!black}\LatoBlack #1}}
% URL mostrata come testo
\def\urlx#1{{\FiraCode\small\color{primary!70!black}\url{#1}}}

% \tasto{Ctrl}: un tasto della tastiera; \tasti{Ctrl,Shift,P}: una
% combinazione di tasti da premere insieme
% (le dimensioni minime si azzerano: dentro un nodo di una figura, il tasto
% erediterebbe quelle del nodo che lo contiene)
\newcommand{\tasto}[1]{\tikz[baseline=(k.base)]\node[draw=neutral!55!black,
  fill=neutral!6!white, rounded corners=2pt, line width=0.5pt,
  minimum width=0pt, minimum height=0pt, text width=, align=center,
  inner xsep=3pt, inner ysep=1.6pt, text=black, font=\Lato\scriptsize] (k) {\strut #1};}
\makeatletter
\newcommand{\tasti}[1]{{\let\wt@sep\relax
  \@for\wt@k:=#1\do{\wt@sep\tasto{\wt@k}\def\wt@sep{\,{\Lato\scriptsize+}\,}}}}
% \menu{File,Open Folder...}: un percorso nei menu di un programma
\newcommand{\menu}[1]{{\let\wt@sep\relax\Lato\bfseries\color{primary!65!black}%
  \@for\wt@k:=#1\do{\wt@sep\mbox{\wt@k}\def\wt@sep{\,{\color{neutral!55!black}›}\,}}}}
\makeatother
% \pulsante{Install}: un pulsante di un'interfaccia grafica
\newcommand{\pulsante}[1]{{\setlength{\fboxsep}{1.6pt}%
  \colorbox{dmaccent}{\Lato\scriptsize\color{white}\strut\,#1\,}}}

% Elenco di passi numerati (installazione, procedure): pallini con il numero
\NewDocumentCommand{\wtpassonum}{m}{\tikz[baseline=(n.base)]\node[circle,
  fill=primary!80!black, text=white, font=\Lato\footnotesize\bfseries,
  inner sep=0pt, minimum size=13pt] (n) {#1};}
\newlist{passi}{enumerate}{1}
\setlist[passi]{label=\wtpassonum{\arabic*}, leftmargin=24pt, labelsep=6pt,
  itemsep=3pt, topsep=4pt}

% ============================================================
%  RIQUADRO "IN SINTESI"
% ============================================================
% Chiude le parti piu' dense riassumendo il ragionamento (non una regola
% da imparare a memoria, ma il filo logico da ricordare).
\newtcolorbox{sintesibox}{enhanced, breakable,
  colback=primary!4!white, colframe=primary!70!black, boxrule=0pt,
  leftrule=3pt, arc=0pt, outer arc=0pt, left=8pt, right=8pt, top=5pt,
  bottom=5pt, before skip=0.4cm, after skip=0.5cm, fontupper=\small,
  title={\Lato\bfseries\small In sintesi}, coltitle=primary!70!black,
  colbacktitle=primary!4!white, titlerule=0pt, toptitle=3pt, bottomtitle=0pt}
\newcommand{\insintesi}[1]{\begin{sintesibox}#1\end{sintesibox}}

% ============================================================
%  FINESTRA "BROWSER" (pagine web da cui si scaricano i programmi)
%  Uso: \begin{browser}[indirizzo] ... \end{browser}
% ============================================================
\newcommand{\wtbrowserbefore}{0.35cm}
\newcommand{\wtbrowserafter}{0.6cm}
\newtcolorbox{browser}[1][]{%
  enhanced,
  colback=white,
  colframe=neutral!35!white,
  boxrule=0.8pt,
  arc=3pt,
  drop fuzzy shadow,
  before skip=\wtbrowserbefore, after skip=\wtbrowserafter,
  left=8pt, right=8pt, top=6pt, bottom=6pt,
  fonttitle=\footnotesize\Lato,
  colbacktitle=neutral!15!white,
  coltitle=neutral!35!black,
  title={%
    \tikz[baseline=-2.5pt]{%
      \fill[red-gradient-3] (0,0) circle (2.2pt);
      \fill[yellow-gradient-5] (8pt,0) circle (2.2pt);
      \fill[green-gradient-6] (16pt,0) circle (2.2pt);}%
    \hspace{8pt}\footnotesize #1},
  attach boxed title to top left={yshift*=-\tcboxedtitleheight},
  boxed title style={
    sharp corners=downhill, arc=3pt,
    colframe=neutral!35!white, boxrule=0.8pt,
  },
  fontupper=\Lato\small,
}

% ============================================================
%  TERMINALE, per mostrare comandi e output dei programmi
%  Uso: \begin{console}[titolo] ... \end{console}
% ============================================================
% Il contenuto e' VERBATIM: l'output di un programma contiene spesso
% caratteri speciali per LaTeX (_ \ ^ $ % &) e gli a-capo vanno preservati.
% Le righe che cominciano con il prompt "$" (o "PS C:" in PowerShell) sono
% COMANDI digitati dall'utente e sono piu' chiare; il resto e' cio' che
% stampa il computer. Cio' che l'utente digita mentre un programma e' in
% esecuzione (l'input) si racchiude fra << e >> e compare in giallo.
% gcc cita i nomi fra virgolette tipografiche (‘main’): listings non sa
% gestire i caratteri non ASCII e li stamperebbe fuori posto, quindi nel
% sorgente vanno scritti come (*@‘@*) e (*@’@*).
% Come per lstlisting, questo ambiente non puo' stare dentro l'argomento
% di una macro: nei riquadri va usata la forma ambiente (warningbox...).
\definecolor{consolebg}{HTML}{22252E}
\definecolor{consolefg}{HTML}{C9CDD6}
\definecolor{consolecmd}{HTML}{FFFFFF}
\definecolor{consoleinput}{HTML}{F2D46B}
\definecolor{consolebar}{HTML}{31353F}
\newcommand{\wtconsolebefore}{0.35cm}
\newtcblisting{console}[1][Terminale]{%
  listing only,
  enhanced,
  colback=consolebg,
  colframe=consolebg,
  boxrule=0pt,
  arc=3pt,
  drop fuzzy shadow,
  before skip=\wtconsolebefore, after skip=0.6cm,
  left=8pt, right=8pt, top=6pt, bottom=6pt,
  fonttitle=\footnotesize\Lato,
  colbacktitle=consolebar,
  coltitle=consolefg!80!consolebg,
  title={\iterminal\hspace{6pt}\footnotesize #1},
  attach boxed title to top left={yshift*=-\tcboxedtitleheight},
  boxed title style={sharp corners=downhill, arc=3pt, boxrule=0pt,
                     colframe=consolebar},
  listing options={
    language={},
    basicstyle=\linespread{1.1}\footnotesize\FiraCodeNL\color{consolefg},
    backgroundcolor=\color{consolebg},
    keywordstyle={}, commentstyle={}, stringstyle={},
    moredelim={[l][\color{consolecmd}\bfseries]{\$\ }},
    moredelim={[l][\color{consolecmd}\bfseries]{PS\ C:}},
    moredelim={[is][\color{consoleinput}\bfseries]{<<}{>>}},
    numbers=none, frame=none,
    xleftmargin=0pt, xrightmargin=0pt,
    breaklines=true, columns=fixed,
    mathescape=false, escapeinside={(*@}{@*)},
    showtabs=false, tabsize=8,
  },
}

% ============================================================
%  STILI TIKZ CONDIVISI DALLE FIGURE
% ============================================================
\tikzset{
  % nei riquadri delle figure le parole non vanno spezzate con il trattino
  nosillabe/.style={execute at begin node={\hyphenpenalty=10000\exhyphenpenalty=10000}},
  wtnode/.style={nosillabe,
    rounded corners=3pt, draw=primary!65!black, fill=primary!8!white,
    line width=0.8pt, inner sep=5pt, align=center, font=\Lato\small},
  wtnodeacc/.style={
    wtnode, draw=accent!60!black, fill=accent!8!white},
  wtnodewarn/.style={
    wtnode, draw=orange-gradient-6, fill=orange-gradient-1!45!white},
  wtnodeok/.style={
    wtnode, draw=green-gradient-6, fill=green-gradient-1!45!white},
  wtnodeplain/.style={nosillabe,
    rounded corners=2pt, draw=neutral!60!white, fill=neutral!8!white,
    line width=0.7pt, inner sep=4pt, align=center, font=\Lato\footnotesize},
  wtarrow/.style={
    -{Stealth[length=6pt,width=5pt]}, line width=0.9pt, draw=primary!70!black},
  wtarrowacc/.style={
    -{Stealth[length=6pt,width=5pt]}, line width=0.9pt, draw=accent!70!black},
  wtdash/.style={
    dashed, line width=0.7pt, draw=neutral!50!white},
  wtlab/.style={
    font=\Lato\scriptsize, text=neutral!15!black, inner sep=2.5pt,
    fill=white, align=center},
  wtbrace/.style={
    decorate, decoration={brace, amplitude=4pt, mirror}, draw=neutral!40!black,
    line width=0.7pt},
}

% Figura centrata con didascalia, ridotta alla larghezza indicata
\newcommand{\wtfig}[3][0.95]{%
  \begin{figure}[H]
    \centering
    \resizebox{#1\linewidth}{!}{#2}
    \caption{#3}
  \end{figure}}

% ============================================================
%  IL CODICE NON SI SPEZZA FRA DUE PAGINE
% ============================================================
% Ogni lstlisting viene racchiuso in una minipage. Una minipage e' una
% scatola indivisibile: se il listato non entra nello spazio rimasto in
% fondo alla pagina, passa per intero alla pagina successiva.
\AddToHook{env/lstlisting/before}{\wtattacca\noindent\begin{minipage}{\linewidth}}
\AddToHook{env/lstlisting/after}{\end{minipage}\par}

% ============================================================
%  CODICE E RISULTATO AFFIANCATI
% ============================================================
% Uso:
%   \begin{affianco}[0.55]      % frazione della riga occupata dal codice
%     \begin{lstlisting}...\end{lstlisting}   (uno o piu' listati)
%   \risultato
%     \begin{console} ... \end{console}   (oppure uno schema)
%   \end{affianco}
% Le due colonne sono allineate in alto e formano un blocco unico, che non
% viene mai diviso fra due pagine. Con il secondo argomento facoltativo [c]
% le colonne sono centrate in verticale l'una rispetto all'altra.
\newlength{\wtcodew}
\newlength{\wtresw}
\newlength{\wtgap}
\setlength{\wtgap}{0.5cm}
\NewDocumentEnvironment{affianco}{O{0.56} O{t}}{%
  \par\addvspace{6pt}\noindent
  \setlength{\wtcodew}{#1\linewidth}%
  \setlength{\wtresw}{\linewidth-\wtcodew-\wtgap}%
  \wtvscmargini{0pt}{0pt}%
  \lstset{aboveskip=0pt, belowskip=5pt}%
  \renewcommand{\wtbrowserbefore}{5pt}%
  \renewcommand{\wtbrowserafter}{6pt}%
  \renewcommand{\wtconsolebefore}{5pt}%
  \if c#2%
    \def\risultato{\end{minipage}\hspace{\wtgap}%
      \begin{minipage}[c]{\wtresw}}%
    \begin{minipage}[c]{\wtcodew}%
  \else
    \def\risultato{\end{minipage}\hspace{\wtgap}%
      \begin{minipage}[t]{\wtresw}\vspace{0pt}}%
    \begin{minipage}[t]{\wtcodew}\vspace{0pt}%
  \fi
}{%
  \end{minipage}\par\addvspace{10pt}}

% Listati un po' piu' piccoli, per le righe lunghe dentro una colonna stretta
\newcommand{\codicepiccolo}{%
  \lstset{basicstyle=\linespread{1.12}\wtvsccodefontsmall,
          numberstyle=\fontsize{5}{6}\selectfont\FiraCodeNL\color{vsclinenum}}}

% Didascalia breve sotto un risultato
\newcommand{\wtnota}[1]{%
  {\par\centering\Lato\footnotesize\color{neutral!30!black}#1\par}\vspace{4pt}}

% ============================================================
%  SCHEMA ACCANTO AL CODICE
% ============================================================
% Nella colonna di destra di "affianco":
%   \begin{immagine}[didascalia facoltativa]
%     \begin{tikzpicture} ... \end{tikzpicture}
%   \end{immagine}
% Lo schema viene centrato nella colonna e, se e' piu' largo, ridotto alla
% sua larghezza (nel log compare "Schema ridotto", per accorgersene).
\newsavebox{\wtimgbox}
\NewDocumentEnvironment{immagine}{o}
  {\begin{lrbox}{\wtimgbox}}
  {\end{lrbox}%
   \par\noindent\makebox[\linewidth]{%
     \ifdim\wd\wtimgbox>\linewidth
       \typeout{Schema ridotto: \the\wd\wtimgbox\space -> \the\linewidth
                (riga \the\inputlineno)}%
       \resizebox{\linewidth}{!}{\usebox{\wtimgbox}}%
     \else
       \usebox{\wtimgbox}%
     \fi}\par
   \IfValueT{#1}{\vspace{4pt}\wtnota{#1}}}

% ============================================================
%  SCHEMI DELLA MEMORIA (variabili, puntatori, array, struct)
% ============================================================
% Livelli per i riquadri annidati (per esempio le chiamate di funzione
% sulla pila): il piu' esterno va disegnato sotto agli altri
\pgfdeclarelayer{liv1}
\pgfdeclarelayer{liv2}
\pgfdeclarelayer{liv3}
\pgfsetlayers{liv1,liv2,liv3,background,main}

\colorlet{mref}{accent!70!black}
\colorlet{mok}{green-gradient-6!85!black}
\colorlet{merr}{red-gradient-6!90!black}
\tikzset{
  % casella di una variabile: \node[mvar] (x) {42};
  mvar/.style={draw=neutral!55!black, fill=white, line width=0.7pt,
    minimum height=0.5cm, minimum width=1.4cm, inner sep=2pt,
    font=\FiraCodeNL\scriptsize},
  % nome della variabile, a sinistra della casella
  mnome/.style={font=\FiraCodeNL\scriptsize, text=primary!60!black,
    inner sep=2pt},
  % indirizzo di memoria, accanto a una casella
  mindir/.style={font=\FiraCodeNL\tiny, text=neutral!45!black, inner sep=2pt},
  % intestazione e righe di una struttura (vedi \oggetto)
  mtesta/.style={draw=primary!65!black, fill=primary!16!white,
    line width=0.7pt, minimum height=0.46cm, inner sep=2pt,
    font=\Lato\scriptsize\bfseries, text=primary!35!black},
  mriga/.style={draw=primary!65!black, fill=primary!3!white,
    line width=0.7pt, minimum height=0.46cm, inner sep=3pt,
    font=\FiraCodeNL\scriptsize, align=left},
  % freccia di un puntatore e pallino da cui parte
  mrif/.style={-{Stealth[length=5pt,width=4pt]}, line width=0.85pt,
    draw=mref},
  mpunto/.style={circle, fill=mref, inner sep=1.4pt},
  % freccia tratteggiata: una copia, un percorso, un "va a leggere"
  mcerca/.style={-{Stealth[length=5pt,width=4pt]}, line width=0.8pt,
    draw=green-gradient-6!80!black, dashed},
  % nota esplicativa nello schema
  mnota/.style={font=\Lato\scriptsize, text=neutral!20!black, align=left,
    inner sep=2pt},
  % codice dentro uno schema
  mcod/.style={font=\FiraCodeNL\scriptsize, inner sep=2pt},
  % riquadro di una funzione
  mfun/.style={draw=accent!60!black, fill=accent!8!white, line width=0.7pt,
    rounded corners=6pt, minimum height=0.5cm, inner sep=3pt,
    font=\FiraCodeNL\scriptsize},
  % riquadri annidati: \scopo{livello}{nome}{etichetta}{(nodo1) (nodo2) ...}
  mscopo1/.style={draw=neutral!50!white, fill=neutral!6!white,
    line width=0.7pt, rounded corners=5pt, inner sep=7pt},
  mscopo2/.style={draw=primary!45!white, fill=primary!6!white,
    line width=0.7pt, rounded corners=5pt, inner sep=7pt},
  mscopo3/.style={draw=accent!45!white, fill=accent!7!white,
    line width=0.7pt, rounded corners=5pt, inner sep=7pt},
  mscopolab/.style={font=\Lato\scriptsize\bfseries, text=neutral!25!black,
    inner sep=1.5pt},
}

% \oggetto[stile della testata]{nome}{(posizione)}{larghezza}{titolo}
%         {campo/valore, campo/valore, ...}
% Disegna una struttura (struct): crea i nodi "nome" (testata), "nome-1",
% "nome-2", ... (righe) e "nome-box" (tutta la struttura). La posizione e'
% quella del centro della testata, salvo un'ancora diversa nel primo
% argomento (per esempio anchor=north west). La larghezza indicata e'
% quella minima: se una riga e' piu' lunga, la struttura si allarga. Un
% valore che contiene una virgola va tra graffe: sep/{","}.
\newlength{\mchiavew}
\setlength{\mchiavew}{1.25cm}
\newcommand{\mchiave}[1]{\setlength{\mchiavew}{#1}}
\newlength{\mobjw}
\newlength{\mkeyw}
\newlength{\mrigaw}
\newsavebox{\wtimgmisura}
\NewDocumentCommand{\oggetto}{O{} m m m m m}{%
  \global\mobjw=#4\relax
  \global\mkeyw=\mchiavew
  \ifx\relax#6\relax\else
    \foreach \mk/\mv in {#6} {%
      \pgfinterruptpicture
        \global\setbox\wtimgmisura\hbox{\FiraCodeNL\scriptsize\mk}%
      \endpgfinterruptpicture
      \ifdim\dimexpr\wd\wtimgmisura+6pt\relax>\mkeyw
        \global\mkeyw=\dimexpr\wd\wtimgmisura+6pt\relax
      \fi}%
    \foreach \mk/\mv in {#6} {%
      \pgfinterruptpicture
        \global\setbox\wtimgmisura\hbox{\FiraCodeNL\scriptsize
          \makebox[\mkeyw][l]{\mk}\mv}%
      \endpgfinterruptpicture
      \setlength{\mrigaw}{\wd\wtimgmisura}%
      \ifdim\dimexpr\mrigaw+8pt\relax>\mobjw
        \global\mobjw=\dimexpr\mrigaw+8pt\relax
      \fi}%
  \fi
  \node[mtesta, minimum width=\mobjw, #1] (#2) at #3 {#5};
  \xdef\wtprev{#2}%
  \ifx\relax#6\relax
    \node[mriga, minimum width=\mobjw, anchor=north] (#2-1)
      at ([yshift=0.7pt]\wtprev.south) {\phantom{x}};
    \xdef\wtprev{#2-1}%
  \else
    \foreach \mk/\mv [count=\mi] in {#6} {%
      \node[mriga, minimum width=\mobjw, text width=\dimexpr\mobjw-6pt\relax,
            anchor=north] (#2-\mi) at ([yshift=0.7pt]\wtprev.south)
        {\makebox[\mkeyw][l]{\color{primary!55!black}\mk}\mv};
      \xdef\wtprev{#2-\mi}%
    }%
  \fi
  \node[fit=(#2)(\wtprev), inner sep=0pt] (#2-box) {};}

% \variabile[stile]{nome del nodo}{(posizione)}{nome mostrato}{valore}
% Il nome compare a sinistra della casella (nodo "nome-nome"). Per un
% puntatore si lascia il valore vuoto e si usa \riferimento.
\NewDocumentCommand{\variabile}{O{} m m m m}{%
  \node[mvar, #1] (#2) at #3 {#5};
  \node[mnome, anchor=east] (#2-nome) at ([xshift=-2pt]#2.west) {#4};}

% \riferimento[stile del percorso]{nodo di partenza}{percorso fino alla meta}
% es.: \riferimento{p}{-- (x.west)}  oppure  \riferimento{p}{to[bend left] (x)}
\NewDocumentCommand{\riferimento}{O{} m m}{%
  \draw[mrif, #1] (#2.center) #3;
  \node[mpunto] at (#2.center) {};}

% \scopo[stile]{livello 1-3}{nome}{etichetta}{(nodo) (nodo) ...}
% Il riquadro racchiude i nodi indicati e lascia in alto lo spazio per
% l'etichetta (per esempio il nome della funzione di cui e' la memoria).
\NewDocumentCommand{\scopo}{O{} m m m m}{%
  \node[fit={#5}, inner sep=0pt] (#3-in) {};
  \node[mscopolab, anchor=south west] (#3-lab)
    at ([yshift=1pt]#3-in.north west) {#4};
  \begin{pgfonlayer}{liv#2}
    \node[mscopo#2, fit=(#3-in)(#3-lab), #1] (#3) {};
  \end{pgfonlayer}}

% \celle[stile]{nome}{(posizione)}{v0, v1, ...}: le caselle consecutive di
% un array, con l'indice sotto ciascuna. La posizione e' il bordo sinistro
% della prima casella; crea i nodi nome-0, nome-1, ... e nome-box. Con
% \mvuota si disegna una casella senza valore, con \noindici prima della
% chiamata si omettono gli indici.
\tikzset{
  mcella/.style={draw=primary!65!black, fill=primary!3!white, line width=0.7pt,
    minimum height=0.5cm, minimum width=0.75cm, inner xsep=4pt,
    font=\FiraCodeNL\scriptsize},
  mindice/.style={font=\FiraCodeNL\tiny, text=neutral!45!black, inner sep=1pt},
}
\newcommand{\mvuota}{{\Lato\itshape\color{neutral!50!black}?}}
\newif\ifmindici \mindicitrue
\newcommand{\noindici}{\mindicifalse}
\NewDocumentCommand{\celle}{O{} m m m}{%
  \coordinate (#2-inizio) at #3;
  \xdef\wtprev{#2-inizio}%
  \foreach \v [count=\i from 0] in {#4} {%
    \node[mcella, anchor=west, #1] (#2-\i)
      at ([xshift=-0.7pt]\wtprev.east |- #2-inizio) {\v};
    \ifmindici
      \node[mindice, anchor=north] at ([yshift=-1pt]#2-\i.south) {\i};
    \fi
    \xdef\wtprev{#2-\i}%
  }
  \node[fit=(#2-0)(\wtprev), inner sep=0pt] (#2-box) {};}

% \memoria[stile]{nome}{(posizione)}{indirizzo/valore/etichetta, ...}: una
% colonna di celle di memoria, dall'alto verso il basso, con l'indirizzo a
% sinistra e (facoltativa) un'etichetta a destra. La posizione e' il centro
% della prima cella; crea i nodi nome-1, nome-2, ... e nome-box.
\tikzset{
  mcellam/.style={draw=primary!65!black, fill=primary!3!white, line width=0.7pt,
    minimum height=0.48cm, minimum width=1.7cm, inner sep=2pt,
    font=\FiraCodeNL\scriptsize},
}
\NewDocumentCommand{\memoria}{O{} m m m}{%
  \coordinate (#2-inizio) at #3;
  \foreach \ind/\val/\etic [count=\i] in {#4} {%
    \node[mcellam, #1] (#2-\i) at ($(#2-inizio)+(0,-\i*0.48cm+0.48cm-\i*0.7pt+0.7pt)$) {\val};
    \node[mindir, anchor=east] at ([xshift=-2pt]#2-\i.west) {\ind};
    \node[mnome, anchor=west] at ([xshift=3pt]#2-\i.east) {\etic};
    \xdef\wtultima{#2-\i}%
  }
  \node[fit=(#2-1)(\wtultima), inner sep=0pt] (#2-box) {};}

% ─── Liste collegate ─────────────────────────────────────────
% \nodo[stile]{nome}{(posizione)}{dato}: un nodo di una lista, con il campo
% dei dati e il campo next affiancati. La posizione e' il bordo sinistro;
% crea i nodi nome-dato, nome-next e nome (tutto il nodo). \nodonull{nome}
% barra il campo next: il puntatore vale NULL. \collega{da}{a} disegna la
% freccia dal campo next del primo nodo al secondo.
\tikzset{
  mlnodo/.style={draw=primary!65!black, fill=primary!3!white, line width=0.7pt,
    minimum height=0.55cm, inner sep=2pt, font=\FiraCodeNL\scriptsize},
  mlnuovo/.style={draw=accent!65!black, fill=accent!10!white},
  mlvia/.style={draw=red-gradient-6!70!black, fill=red-gradient-1!25!white},
}
\NewDocumentCommand{\nodo}{O{} m m m}{%
  \node[mlnodo, minimum width=0.95cm, anchor=west, #1] (#2-dato) at #3 {#4};
  \node[mlnodo, minimum width=0.55cm, anchor=west, #1] (#2-next)
    at ([xshift=-0.7pt]#2-dato.east) {};
  \node[fit=(#2-dato)(#2-next), inner sep=0pt] (#2) {};}
\newcommand{\nodonull}[1]{%
  \draw[line width=0.7pt, draw=primary!65!black]
    ([shift={(1.5pt,1.5pt)}]#1-next.south west) -- ([shift={(-1.5pt,-1.5pt)}]#1-next.north east);}
\NewDocumentCommand{\collega}{O{} m m}{%
  \draw[mrif, #1] (#2-next.center) -- (#3.west);
  \node[mpunto] at (#2-next.center) {};}

% Segni di riuscita e di errore negli schemi
\newcommand{\mok}{{\color{mok}\itick}}
\newcommand{\merr}{{\color{merr}\icross}}

% ============================================================
%  DIAGRAMMI DI FLUSSO
% ============================================================
% Forme standard: inizio e fine (ovale), istruzione (rettangolo), input e
% output (parallelogramma), decisione (rombo), sottoprogramma (rettangolo
% con le bande laterali).
\tikzset{
  fcterm/.style={nosillabe, draw=primary!65!black, fill=primary!18!white, line width=0.8pt,
    rounded corners=0.29cm, minimum height=0.58cm, minimum width=1.7cm,
    inner xsep=8pt, font=\Lato\scriptsize\bfseries, text=primary!35!black},
  fcproc/.style={nosillabe, draw=primary!65!black, fill=primary!4!white, line width=0.8pt,
    minimum height=0.58cm, minimum width=2.1cm, inner sep=4pt, align=center,
    font=\Lato\scriptsize},
  fcio/.style={nosillabe, trapezium, trapezium left angle=72, trapezium right angle=108,
    trapezium stretches body, draw=green-gradient-6!80!black,
    fill=green-gradient-1!40!white, line width=0.8pt, minimum height=0.58cm,
    minimum width=2.1cm, inner sep=3pt, align=center, font=\Lato\scriptsize},
  fcdec/.style={nosillabe, diamond, aspect=2.1, draw=accent!60!black, fill=accent!7!white,
    line width=0.8pt, inner sep=1.5pt, align=center, font=\Lato\scriptsize},
  fcsub/.style={fcproc, fill=orange-gradient-1!30!white,
    draw=orange-gradient-6!80!black, inner xsep=9pt,
    path picture={\draw[orange-gradient-6!80!black, line width=0.6pt]
      ([xshift=4pt]path picture bounding box.north west) --
      ([xshift=4pt]path picture bounding box.south west)
      ([xshift=-4pt]path picture bounding box.north east) --
      ([xshift=-4pt]path picture bounding box.south east);}},
  fcfreccia/.style={-{Stealth[length=5pt,width=4.2pt]}, line width=0.8pt,
    draw=neutral!30!black},
  fclinea/.style={line width=0.8pt, draw=neutral!30!black},
  fcsi/.style={font=\Lato\tiny\bfseries, text=green-gradient-6!55!black,
    inner sep=1.5pt},
  fcno/.style={font=\Lato\tiny\bfseries, text=red-gradient-6!75!black,
    inner sep=1.5pt},
  fccod/.style={font=\FiraCodeNL\scriptsize},
  % punto in cui due strade si ricongiungono (o in cui rientra un ciclo)
  fcgiunzione/.style={circle, draw=neutral!30!black, fill=white, line width=0.8pt,
    inner sep=0pt, minimum size=5.5pt},
}

% ============================================================
%  FILE SYSTEM: CARTELLE E FILE
% ============================================================
\tikzset{
  fsdir/.style={inner sep=1.5pt, font=\Lato\scriptsize, text=primary!45!black,
    execute at begin node={{\color{orange-gradient-5!85!black}\ifolder}\,}},
  fsfile/.style={inner sep=1.5pt, font=\Lato\scriptsize, text=neutral!20!black,
    execute at begin node={{\color{neutral!55!black}\ifile}\,}},
  fsqui/.style={fsdir, fill=yellow-gradient-1!70!white, rounded corners=2pt},
  fsramo/.style={draw=neutral!50!black, line width=0.55pt},
}

% ============================================================
%  INTERFACCIA DI VS CODE (tema scuro predefinito "Dark Modern")
% ============================================================
% Colori delle zone della finestra, per le figure che la riproducono
\definecolor{dmbg}{HTML}{1F1F1F}       % editor
\definecolor{dmbar}{HTML}{181818}      % barre: attivita', laterale, stato
\definecolor{dmborder}{HTML}{2B2B2B}   % bordi fra le zone
\definecolor{dmfg}{HTML}{CCCCCC}       % testo
\definecolor{dmdim}{HTML}{9D9D9D}      % testo secondario, icone inattive
\definecolor{dmaccent}{HTML}{0078D4}   % pulsanti, indicatore remoto
\definecolor{dminput}{HTML}{313131}    % caselle di testo
\definecolor{dmsel}{HTML}{2A2D2E}      % voce evidenziata in un elenco
\definecolor{dmyellow}{HTML}{CCA700}   % avvisi

% \vscico[colore]{nome}: le icone della barra delle attivita' di VS Code,
% disegnate in un quadrato 10x10. Nomi: esplora, cerca, git, debug,
% estensioni, remoto, account, impostazioni, gioca.
\makeatletter
\def\vsc@ico@esplora{%
  \draw[line join=round] (2.6,1) rectangle (8.6,7.4);
  \draw[line join=round] (1.2,2.6) -- (1.2,9) -- (6.2,9) -- (7.4,7.8);}
\def\vsc@ico@cerca{%
  \draw (5.9,6.1) circle (2.7);
  \draw[line width=1.1pt] (3.9,4.1) -- (1.2,1.4);}
\def\vsc@ico@git{%
  \draw (3,8.4) circle (1.2); \draw (3,1.6) circle (1.2);
  \draw (7.4,6.6) circle (1.2);
  \draw (3,7.2) -- (3,2.8);
  \draw (7.4,5.4) .. controls (7.4,3.6) and (3,4.6) .. (3,2.8);}
\def\vsc@ico@debug{%
  \draw[line join=round] (1.4,1.2) -- (1.4,8.8) -- (7.4,5) -- cycle;
  \fill (7.6,2.2) circle (1.5);}
\def\vsc@ico@estensioni{%
  \draw (1,1) rectangle (4.6,4.6); \draw (4.6,1) rectangle (8.2,4.6);
  \draw (1,4.6) rectangle (4.6,8.2);
  \draw[rotate around={45:(7.4,7.4)}] (5.6,5.6) rectangle (9.2,9.2);}
\def\vsc@ico@remoto{%
  \draw[line cap=round, line join=round] (1,6.4) -- (3.6,3.8) -- (1,1.2);
  \draw[line cap=round, line join=round] (9,8.8) -- (6.4,6.2) -- (9,3.6);}
\def\vsc@ico@account{%
  \draw (5,6.6) circle (2.2);
  \draw (1.2,0.8) .. controls (1.6,4.2) and (8.4,4.2) .. (8.8,0.8);}
\def\vsc@ico@impostazioni{%
  \draw (5,5) circle (1.6);
  \draw (5,5) circle (3.4);
  \foreach \ang in {0,45,...,315} \draw[line width=1.2pt] (\ang:3.4)++(5,5) -- ++(\ang:1.1);}
\def\vsc@ico@gioca{%
  \fill[line join=round] (2,1) -- (2,9) -- (8.6,5) -- cycle;}
\NewDocumentCommand{\vscico}{O{dmdim} m}{%
  \tikz[baseline=2.4pt, x=1.15pt, y=1.15pt, line width=0.75pt, draw=#1, fill=#1]{%
    \useasboundingbox (0,0) rectangle (10,10);
    \csname vsc@ico@#2\endcsname}}
\makeatother

% Colori del prompt di Ubuntu nel terminale di VS Code (utente@computer in
% verde, cartella in blu)
\definecolor{termgreen}{HTML}{23D18B}
\definecolor{termblue}{HTML}{3B8EEA}

% \vscodelogo[dimensione]: il logo di Visual Studio Code, semplificato
\NewDocumentCommand{\vscodelogo}{O{0.5cm}}{%
  \tikz[baseline=0pt, x=#1, y=#1]{%
    \fill[dmaccent!85!white] (0.70,0.98) -- (0.95,0.86) -- (0.95,0.14) -- (0.70,0.02)
      -- (0.31,0.37) -- (0.13,0.23) -- (0.04,0.28) -- (0.04,0.72) -- (0.13,0.77)
      -- (0.31,0.63) -- cycle;
    \fill[white] (0.70,0.73) -- (0.70,0.27) -- (0.43,0.5) -- cycle;
    \fill[white] (0.13,0.62) -- (0.23,0.5) -- (0.13,0.38) -- cycle;}}

% Icone della barra di stato: errori (cerchio con la croce) e avvisi
% (triangolo con il punto esclamativo)
\NewDocumentCommand{\icoerr}{O{dmfg}}{\tikz[baseline=-0.75ex]{%
  \draw[#1, line width=0.5pt] (0,0) circle (2.7pt);
  \draw[#1, line width=0.5pt] (-1.2pt,-1.2pt) -- (1.2pt,1.2pt) (-1.2pt,1.2pt) -- (1.2pt,-1.2pt);}}
\NewDocumentCommand{\icowarn}{O{dmfg}}{\tikz[baseline=-0.75ex]{%
  \draw[#1, line width=0.5pt, line join=round] (-3pt,-2.5pt) -- (3pt,-2.5pt) -- (0,3pt) -- cycle;
  \draw[#1, line width=0.5pt] (0,-0.2pt) -- (0,1.4pt); \fill[#1] (0,-1.3pt) circle (0.35pt);}}

% Sottolineatura ondulata degli errori nell'editor
\usetikzlibrary{decorations.pathmorphing}
\tikzset{ondina/.style={decorate, decoration={snake, amplitude=0.55pt,
  segment length=2.6pt}, line width=0.5pt}}

% \wtcasellina{0|1}: casella di spunta di un programma di installazione
\newcommand{\wtcasellina}[1]{\tikz[baseline=-2.8pt]{%
  \ifnum#1=1
    \fill[rounded corners=1pt, blue!70!black] (-3.2pt,-3.2pt) rectangle (3.2pt,3.2pt);
    \draw[white, line width=1.1pt] (-1.9pt,-0.2pt) -- (-0.5pt,-1.7pt) -- (2pt,1.8pt);
  \else
    \draw[rounded corners=1pt, neutral!60!black, line width=0.6pt]
      (-3.2pt,-3.2pt) rectangle (3.2pt,3.2pt);
  \fi}}
% Puntatore del mouse, per indicare dove fare clic
\newcommand{\wtpuntatore}{\tikz[baseline=0pt, scale=0.55]{%
  \filldraw[fill=white, draw=black, line width=0.35pt, line join=round]
    (0,0) -- (0,-0.62) -- (0.15,-0.48) -- (0.26,-0.72) -- (0.35,-0.68)
    -- (0.24,-0.44) -- (0.44,-0.44) -- cycle;}}

% ============================================================
%  ICONE DELLA BARRA DI DEBUG DI VS CODE
% ============================================================
% \wtdbgicon{nome} su sfondo chiaro, \wtdbgiconscura{nome} su sfondo scuro.
% Nomi: maniglia, resume, stepover, stepinto, stepout, restart, disconnect.
\definecolor{dbgblu}{HTML}{007ACC}
\definecolor{dbgverde}{HTML}{388A34}
\definecolor{dbgrosso}{HTML}{A1260D}
\definecolor{dbggrigio}{HTML}{8A8A8A}
\definecolor{dbgbluscuro}{HTML}{75BEFF}
\definecolor{dbgverdescuro}{HTML}{89D185}
\definecolor{dbgrossoscuro}{HTML}{F48771}
\definecolor{dbggrigioscuro}{HTML}{8A8A8A}
\makeatletter
\tikzset{dbgfreccia/.style={-{Stealth[length=3pt,width=3.6pt]}}}
\def\wtdbg@colore@maniglia{grigio}
\def\wtdbg@colore@resume{blu}
\def\wtdbg@colore@stepover{blu}
\def\wtdbg@colore@stepinto{blu}
\def\wtdbg@colore@stepout{blu}
\def\wtdbg@colore@restart{verde}
\def\wtdbg@colore@disconnect{rosso}
\def\wtdbg@maniglia{\foreach \x in {3.5,6.5} \foreach \y in {2,5,8} \fill (\x,\y) circle (0.9);}
\def\wtdbg@resume{\draw (1.5,1) -- (1.5,9);
  \draw[line join=round] (4,1.3) -- (4,8.7) -- (9.6,5) -- cycle;}
\def\wtdbg@stepover{\draw[dbgfreccia] (1.2,4.6) arc[start angle=180, end angle=8, radius=3.8];
  \fill (5,1.5) circle (1.5);}
\def\wtdbg@stepinto{\draw[dbgfreccia] (5,10) -- (5,4.2); \fill (5,1.5) circle (1.5);}
\def\wtdbg@stepout{\draw[dbgfreccia] (5,4.4) -- (5,10.2); \fill (5,1.5) circle (1.5);}
\def\wtdbg@restart{\draw[-{Stealth[length=3pt,width=3.6pt]}]
  (2.2,7.6) arc[start angle=135, end angle=-180, radius=3.9];}
\def\wtdbg@disconnect{%
  \begin{scope}[rotate around={45:(5,5)}]
    \draw[rounded corners=0.8] (3.4,3.6) rectangle (6.4,6.4);
    \draw (6.4,4.4) -- (8.4,4.4) (6.4,5.6) -- (8.4,5.6) (0.8,5) -- (3.4,5);
  \end{scope}
  \draw (1.6,8.4) -- (8.4,1.6);}
\newcommand{\wtdbg@icona}[2]{%
  \edef\wtdbg@c{dbg\csname wtdbg@colore@#1\endcsname#2}%
  \tikz[baseline=0pt, x=0.85pt, y=0.85pt, line width=0.85pt, line cap=round,
        draw=\wtdbg@c, fill=\wtdbg@c]{%
    \useasboundingbox (0,0) rectangle (10,10);
    \csname wtdbg@#1\endcsname}}
\newcommand{\wtdbgicon}[1]{\wtdbg@icona{#1}{}}
\newcommand{\wtdbgiconscura}[1]{\wtdbg@icona{#1}{scuro}}
\makeatother

% ============================================================
%  LE SPIEGAZIONI NON SI SPEZZANO FRA DUE PAGINE
% ============================================================
% I riquadri della classe (info, warning, error, definizione) sono
% dichiarati "breakable": qui la chiave "breakable" di tcolorbox viene
% ridefinita in modo che equivalga a "unbreakable": un riquadro che non
% entra passa intero alla pagina successiva, come i listati.
\tcbset{breakable/.style={breakable@false}}

% ============================================================
%  Un concetto non si spezza fra due pagine
% ============================================================
% Il testo che introduce un listato, una figura, una lista o un risultato
% ("... come segue:") resta attaccato a cio' che introduce: se non entrano
% insieme in fondo alla pagina, passano insieme alla successiva.
\raggedbottom
% Un paragrafo non si divide mai fra due pagine
\interlinepenalty=10000
% Un codice inline lungo in fondo alla riga non esce dal margine
\emergencystretch=3em
\makeatletter
% Niente salto pagina fra la frase introduttiva e la lista o il blocco
% centrato (tabelle, figure non flottanti) che la segue
\@beginparpenalty=10000
% Una lista resta tutta sulla stessa pagina
\@itempenalty=10000
\makeatother
\newcommand{\wtattacca}{\par\nopagebreak}
% (per i listati \wtattacca e' gia' nell'hook che li chiude in una minipage)
\AddToHook{env/figure/before}{\wtattacca}
\AddToHook{env/table/before}{\wtattacca}
\AddToHook{env/console/before}{\wtattacca}
\AddToHook{env/browser/before}{\wtattacca}
\AddToHook{env/affianco/before}{\wtattacca}
% Un titolo non resta da solo in fondo alla pagina, e ogni "blocco" (da un
% titolo di sezione o sottosezione al titolo successivo) che sta in una
% pagina non viene spezzato: la posizione di inizio e di fine di ogni blocco
% viene salvata nel file .aux (zref-savepos) e, alla compilazione
% successiva, prima del titolo si chiede tanto spazio quanto e' alto il
% blocco. Se non c'e', il blocco intero passa alla pagina dopo. Servono
% quindi diverse compilazioni perche' le posizioni si stabilizzino.
\usepackage{zref-savepos,zref-abspage}
\newcount\wtblocco
\makeatletter
\zref@addprop{savepos}{abspage}
\newcommand{\wtfineblocco}{%
  \ifnum\wtblocco>\z@ \zsavepos{wtb\the\wtblocco e}\fi}
\newcommand{\wtiniziablocco}{%
  \global\advance\wtblocco\@ne \zsavepos{wtb\the\wtblocco s}}
% \wtspazio{minimo}: Needspace pari all'altezza del prossimo blocco, se sta
% in una pagina, altrimenti pari al minimo. L'altezza misurata quando il
% blocco stava tutto in una pagina viene ricordata (\wtmem nel file .aux):
% cosi' un blocco gia' spostato non torna indietro alla compilazione dopo.
\newdimen\wtaltezza
\newcommand{\wtmem}[2]{\expandafter\gdef\csname wtmem@#1\endcsname{#2}}
% \wtcap{n}: il blocco n apre un capitolo, quindi il blocco n-1 chiude il
% precedente e la sua "fine" viene registrata nella pagina del nuovo capitolo
\newcommand{\wtcap}[1]{\expandafter\gdef\csname wtcap@#1\endcsname{}}
% \wtsegue{n}: il blocco n contiene solo il titolo di una sezione, seguito
% subito da una sottosezione: il titolo viaggia insieme al blocco successivo.
\newcommand{\wtsegue}[1]{\expandafter\gdef\csname wtsegue@#1\endcsname{}}
\newcommand{\wtspazio}[1]{%
  \begingroup
  \edef\wtn{\the\numexpr\wtblocco+1\relax}%
  \edef\wtps{\zref@extractdefault{wtb\wtn s}{abspage}{0}}%
  \edef\wtpe{\zref@extractdefault{wtb\wtn e}{abspage}{0}}%
  \wtaltezza=#1\relax
  \ifcsname wtsegue@\wtn\endcsname
    \edef\wtm{\the\numexpr\wtn+1\relax}%
    \edef\wtqs{\zref@extractdefault{wtb\wtm s}{abspage}{0}}%
    \edef\wtqe{\zref@extractdefault{wtb\wtm e}{abspage}{0}}%
    \ifnum\wtqs>\z@
      \ifnum\wtqe=\wtqs
        \wtaltezza=\dimexpr\zposy{wtb\wtm s}sp-\zposy{wtb\wtm e}sp
                   +4\baselineskip\relax
      \else\ifcsname wtmem@\wtm\endcsname
        \ifdim\csname wtmem@\wtm\endcsname<\maxdimen
          \wtaltezza=\dimexpr\csname wtmem@\wtm\endcsname+4\baselineskip\relax
        \fi
      \fi\fi
    \fi
  \else
  \ifnum\wtps>\z@
    \ifnum\wtpe=\wtps
      \wtaltezza=\dimexpr\zposy{wtb\wtn s}sp-\zposy{wtb\wtn e}sp\relax
      \ifcsname wtmem@\wtn\endcsname
        \ifdim\csname wtmem@\wtn\endcsname>\wtaltezza
          \wtaltezza=\csname wtmem@\wtn\endcsname\relax
        \fi
      \fi
      \immediate\write\@auxout{\string\wtmem{\wtn}{\the\wtaltezza}}%
    \else\ifnum\wtpe=\numexpr\wtps+1\relax
      \ifcsname wtcap@\the\numexpr\wtn+1\relax\endcsname
      \else\ifdim\zposy{wtb\wtn e}sp>\wttesta
        \ifdim\zposy{wtb\wtn s}sp>\wttesta
          \wtaltezza=.96\textheight
          \immediate\write\@auxout{\string\wtmem{\wtn}{\the\wtaltezza}}%
        \fi
      \else\ifcsname wtmem@\wtn\endcsname
        \wtaltezza=\csname wtmem@\wtn\endcsname\relax
        \immediate\write\@auxout{\string\wtmem{\wtn}{\the\wtaltezza}}%
      \else\ifdim\zposy{wtb\wtn s}sp>\wttesta
        \immediate\write\@auxout{\string\wtmem{\wtn}{\the\maxdimen}}%
      \else
        \wtaltezza=\maxdimen
      \fi\fi\fi\fi
    \fi\fi
  \fi
  \fi
  \ifdim\wtaltezza=\maxdimen
    \par\newpage
  \else
    \ifdim\wtaltezza<#1\relax \wtaltezza=#1\relax\fi
    \ifdim\wtaltezza>.97\textheight \wtaltezza=#1\relax\fi
    \Needspace{\dimexpr\wtaltezza+\baselineskip\relax}%
  \fi
  \endgroup}
% posizione verticale oltre la quale un punto e' "in cima alla pagina"
\newdimen\wttesta
\AtBeginDocument{\wttesta=\dimexpr\paperheight-\topmargin-1in-\headheight-\headsep-40pt\relax}
\AddToHook{cmd/chapter/before}{\wtfineblocco\wtiniziablocco\immediate\write\@auxout{\string\wtcap{\the\wtblocco}}}
\AddToHook{cmd/section/before}{\wtfineblocco\wtspazio{22\baselineskip}\wtiniziablocco}
\AddToHook{cmd/subsection/before}{\wtfineblocco
  \if@nobreak\immediate\write\@auxout{\string\wtsegue{\the\wtblocco}}\fi
  \wtspazio{8\baselineskip}\wtiniziablocco}
\AddToHook{enddocument}{\wtfineblocco}
\makeatother
